\documentclass[12pt,a4paper]{article}

\usepackage[spanish,es-nodecimaldot]{babel}
\usepackage[utf8]{inputenc}
\usepackage[T1]{fontenc}
\usepackage{lmodern}
\usepackage{geometry}
\usepackage{microtype}
\usepackage{setspace}
\usepackage{parskip}
\usepackage{enumitem}
\usepackage{xcolor}
\usepackage{listings}
\usepackage{hyperref}
\usepackage{amsmath}
\usepackage{amssymb}
\usepackage{array}
\usepackage{booktabs}
\usepackage{longtable}
\usepackage{tcolorbox}
\usepackage{tikz}

\geometry{margin=2.5cm}
\onehalfspacing
\setlist[itemize]{leftmargin=*,itemsep=0.2em}
\setlist[enumerate]{leftmargin=*,itemsep=0.2em}

\hypersetup{
  colorlinks=true,
  linkcolor=black,
  urlcolor=blue,
  pdftitle={Resumen teorico TP3 - Testing de Software 2024},
  pdfauthor={}
}

\newtcolorbox{definicionbox}[1][]{
  colback=blue!4!white,
  colframe=blue!50!black,
  fonttitle=\bfseries,
  title=Definicion,
  #1
}

\newtcolorbox{ideabox}[1][]{
  colback=green!4!white,
  colframe=green!40!black,
  fonttitle=\bfseries,
  title=Idea clave,
  #1
}

\newtcolorbox{ejemplobox}[1][]{
  colback=orange!4!white,
  colframe=orange!50!black,
  fonttitle=\bfseries,
  title=Ejemplo,
  #1
}

\lstdefinestyle{codestyle}{
  basicstyle=\ttfamily\small,
  frame=single,
  breaklines=true,
  showstringspaces=false
}

\title{\textbf{Resumen Teorico -- Practico 3}\\
Diseno de tests basado en criterios\\
Particionado del Espacio de Entradas\\
\textit{Testing de Software 2024}}
\author{}
\date{}

\begin{document}

\maketitle
\tableofcontents
\newpage

\section*{Introduccion}
\addcontentsline{toc}{section}{Introduccion}

Este resumen sintetiza el material de teoria asociado al TP3:

\begin{itemize}
  \item \textit{Introduction to Software Testing} (Ammann \& Offutt, 2da edicion), Capitulo 6:
  \textbf{Input Space Partitioning}.
  \item \textit{Notas 05: Diseno de tests basado en criterios} (V. Bengolea).
  \item \textit{Notas 06: Particionando el espacio de entrada} (V. Bengolea).
\end{itemize}

El objetivo del resumen es servir de ``hoja de ruta'' para repasar el practico: definiciones,
intuiciones, ejemplos minimos y como se relacionan los conceptos entre si. La presentacion va de
lo general (Model-Driven Test Design) a lo especifico (criterios sobre dominios de entrada como
ECC, PWC, BCC y ACoC).

\section{Marco general: Diseno de tests basado en criterios}

\subsection{El problema y la idea central}

Para casi cualquier programa, el conjunto de entradas posibles es practicamente infinito. No podemos
probar todas las entradas. La pregunta operacional del testing es:

\begin{quote}
\textit{Dado un presupuesto finito de tests, ¿como elegirlos para maximizar la probabilidad de
detectar fallas?}
\end{quote}

El enfoque \textbf{Model-Driven Test Design (MDTD)} responde con una idea simple: \emph{modelar} el
software de alguna forma estructurada y \emph{recorrer} ese modelo siguiendo reglas explicitas. Las
reglas son los \textbf{criterios de cobertura}.

\begin{ideabox}
``La tarea del tester es definir un modelo del software y encontrar formas de recorrerlo.''
\end{ideabox}

\subsection{Requisitos de test y criterios de cobertura}

\begin{definicionbox}
\textbf{Requisito de test (TR).} Un elemento especifico del artefacto software (o de su modelo)
que un test debe satisfacer o cubrir.

\textbf{Criterio de cobertura.} Una coleccion de reglas que impone los requisitos de test sobre el
conjunto de tests.

\textbf{Cobertura.} Dado un conjunto de requisitos de test \(TR\) para un criterio \(C\), un
conjunto de tests \(T\) satisface \(C\) sii
\[
\forall\, tr \in TR,\ \exists\, t \in T \mid t \text{ satisface } tr.
\]
\end{definicionbox}

Un mismo modelo se puede recorrer con criterios distintos, y cada criterio produce su propio
conjunto de TRs. Cambiar el criterio cambia el costo (numero de tests) y la rigurosidad de la
suite.

\subsection{Flujo MDTD}

El flujo recomendado tiene dos niveles de abstraccion:

\begin{enumerate}
  \item \textbf{Nivel de diseno.}
  \begin{enumerate}
    \item Partir del artefacto software (codigo, especificacion, casos de uso, etc.) y construir un
    \textbf{modelo / estructura}.
    \item Aplicar un criterio para derivar \textbf{requisitos de test}.
    \item Refinarlos hasta tener \textbf{especificaciones de test} (que se va a probar y como).
  \end{enumerate}
  \item \textbf{Nivel de implementacion.}
  \begin{enumerate}
    \setcounter{enumii}{3}
    \item Elegir \textbf{valores de inputs} para cada especificacion.
    \item Codificarlos como \textbf{scripts de test} ejecutables.
    \item Correr la suite y obtener \textbf{resultados} (pasa / falla).
  \end{enumerate}
\end{enumerate}

La separacion entre diseno e implementacion es importante: las decisiones de modelado se toman a
nivel de especificacion, y solo despues bajamos al codigo.

\subsection{Cuatro estructuras para modelar software}

El libro distingue cuatro formas de modelar el software, cada una con su propio espacio de
criterios:

\begin{itemize}
  \item \textbf{Espacio de inputs (conjuntos).} Caracterizar el dominio de las entradas y
  particionarlo. \emph{Es el tema central del Capitulo 6 y del TP3.}
  \item \textbf{Grafos.} Modelar flujo de control, flujo de datos, maquinas de estados, etc.
  \item \textbf{Logica.} Cubrir expresiones booleanas (condiciones en codigo, guardas, FNDs, FSMs).
  \item \textbf{Sintaxis.} Cubrir gramaticas (codigo, modelos, formatos de entrada, integracion).
\end{itemize}

\subsection{Cajas coloreadas y MDTD}

La terminologia clasica distingue:
\begin{itemize}
  \item \textbf{Black-box testing}: derivar tests de descripciones externas al codigo
  (especificaciones, requisitos, diseno).
  \item \textbf{White-box testing}: derivar tests de la estructura interna del codigo (sentencias,
  condiciones).
  \item \textbf{Model-based testing}: derivar tests de un modelo abstracto (por ejemplo, un
  diagrama UML).
\end{itemize}
MDTD le resta importancia a esta clasificacion: lo que importa es \emph{que modelo} se usa y
\emph{como} se recorre, no si la fuente es ``externa'' o ``interna''.

\subsection{Requisitos infeasibles y grado de cobertura}

\begin{definicionbox}
\textbf{Requisito de test inviable (infeasible).} Aquel para el cual no existen valores de
entrada que lo satisfagan.
\end{definicionbox}

Tipico ejemplo: una rama de codigo que es \textit{dead code}. La deteccion automatica de requisitos
inviables es \textbf{indecidible} en general. Por eso casi nunca se alcanza el 100\,\% de cobertura
en la practica.

\begin{definicionbox}
\textbf{Grado de cobertura.} Porcentaje de los requisitos de test satisfechos por un conjunto de
tests dado.
\end{definicionbox}

\subsection{Subsumption, conjuntos minimos y minimales}

\begin{definicionbox}
\textbf{Subsumption.} El criterio \(C_1\) \textit{subsume} a \(C_2\) si y solo si \emph{todo}
conjunto de tests que satisface \(C_1\) tambien satisface \(C_2\).
\end{definicionbox}

Es una relacion entre criterios, no entre suites concretas. La intuicion es: ``si paso \(C_1\),
gratis paso \(C_2\)''. Hay dos formas tipicas de que esto suceda:

\begin{itemize}
  \item Los TRs de \(C_1\) son un \emph{superconjunto} de los de \(C_2\).
  \item Existe una correspondencia muchos-a-uno entre TRs de \(C_1\) y de \(C_2\). Si la
  correspondencia es uno-a-uno, los criterios se subsumen mutuamente.
\end{itemize}

\begin{ejemplobox}
\textit{Caramelos.} Si tenemos 6 sabores agrupados en 4 colores, comer un caramelo de cada sabor
implica haber comido al menos uno de cada color (porque cada color esta cubierto por algun sabor).
El criterio ``probar todos los sabores'' \textbf{subsume} al criterio ``probar todos los colores''.

En cambio, \(T_2 = \{\text{limon, pistacho, pera, mandarina}\}\) cumple el criterio de color
(amarillo, verde, blanco, naranja) pero \emph{no} cumple el de sabor (falta damasco y melon). Por
lo tanto el criterio de color \emph{no} subsume al de sabor.
\end{ejemplobox}

\begin{definicionbox}
\textbf{Conjunto de tests minimal.} \(T\) satisface \(TR\) y si se elimina cualquier test, deja
de satisfacerlo. \\
\textbf{Conjunto de tests minimo.} \(T\) satisface \(TR\) y no existe otro \(T'\) mas chico que
tambien lo satisfaga.
\end{definicionbox}

Todo conjunto minimo es minimal, pero no al reves: puede haber varios conjuntos minimales de
tamanos distintos.

\subsection{Generadores, reconocedores y herramientas}

\begin{itemize}
  \item \textbf{Generador}: procedimiento que produce valores que satisfacen un criterio. Es lo
  que tipicamente uno espera, pero suele ser dificil sin herramientas.
  \item \textbf{Reconocedor}: procedimiento que decide si un conjunto de tests dado satisface un
  criterio. Suele ser mas facil.
\end{itemize}

En la industria es comun escribir tests ``a ojo'' y luego \emph{medir} la cobertura con un
reconocedor (por ejemplo, JaCoCo para cobertura de codigo).

\subsection{Por que conviene disenar tests basados en criterios}

\begin{itemize}
  \item Maximiza el beneficio del tiempo invertido en testing: menos tests que son mas efectivos.
  \item Produce suites amplias con minimo solapamiento.
  \item \textbf{Trazabilidad}: cada test esta justificado por uno o varios TRs. Esto facilita el
  testing de regresion y la revision de la suite.
  \item Da una \textbf{regla de parada} explicita: cuantos tests hace falta tener (o cuanta
  cobertura).
\end{itemize}

\subsection{Caracteristicas de un buen criterio de cobertura}

Un buen criterio:
\begin{enumerate}
  \item Permite \emph{calcular los TRs automaticamente} (o al menos de forma sistematica).
  \item Permite \emph{generar valores de test eficientemente}.
  \item Produce suites con \emph{buena capacidad de deteccion} de defectos.
\end{enumerate}

\section{Particionado del Espacio de Entradas}

\subsection{Motivacion}

El \textbf{dominio de entradas} de un programa contiene todas las entradas posibles. Incluso para
programas pequeños es practicamente infinito. La idea del particionado es:

\begin{ideabox}
Dividir el dominio en \textbf{bloques} mediante \textbf{caracteristicas} relevantes para testing y
elegir al menos un valor de cada bloque. Dentro de un mismo bloque, los valores se asumen
``igualmente utiles'' para detectar defectos.
\end{ideabox}

\subsection{Particiones, propiedades formales}

Para un dominio \(D\), una particion \(q\) define un conjunto de bloques
\(B_q = \{b_1, b_2, \ldots, b_n\}\) que debe satisfacer:

\begin{enumerate}
  \item \textbf{Disjointness (no solapamiento).}
  \[
  b_i \cap b_j = \emptyset \quad \forall\, i \neq j,\ b_i, b_j \in B_q.
  \]
  Ningun valor cae en dos bloques a la vez.

  \item \textbf{Completitud.}
  \[
  \bigcup_{b \in B_q} b = D.
  \]
  Todo valor del dominio cae en algun bloque.
\end{enumerate}

\begin{ejemplobox}
Para \texttt{abs(int n)} la caracteristica natural es ``relacion de \(n\) con cero''. Bloques:
\(\{n = 0,\ n > 0,\ n < 0\}\). Es completa y disjunta para \texttt{int}.
\end{ejemplobox}

\subsection{Particiones mal definidas: un caso instructivo}

Un error tipico es elegir bloques que no son ni disjuntos ni completos. Considerese la
caracteristica ``orden de un archivo \(F\)'' con bloques:
\begin{itemize}
  \item \(b_1\): orden creciente,
  \item \(b_2\): orden decreciente,
  \item \(b_3\): orden arbitrario.
\end{itemize}

Si \(F\) tiene \textbf{longitud 1}, el archivo esta simultaneamente en los tres bloques (un
elemento es ``creciente'', ``decreciente'' y ``arbitrario'' a la vez). Hay solapamiento.

\textbf{Solucion}: dividir en dos caracteristicas independientes, cada una con dos bloques:
\(C_1\) = ``orden creciente'' \(\in \{\text{true},\text{false}\}\), \(C_2\) = ``orden
decreciente'' \(\in \{\text{true},\text{false}\}\). Ahora cada particion es disjunta y la
combinacion de ambas describe correctamente las cuatro situaciones posibles.

\begin{ideabox}
Regla practica: \emph{una caracteristica debe modelar una sola propiedad}. Si una particion
parece ``rara'', es mas seguro tener mas caracteristicas con pocos bloques que pocas
caracteristicas con muchos bloques.
\end{ideabox}

\subsection{Modelado del dominio de las entradas: los 5 pasos}

Modelar bien es la mitad del trabajo. El proceso canonico tiene 5 pasos:

\begin{enumerate}
  \item \textbf{Identificar la funcionalidad a testear.}
  \begin{itemize}
    \item Testing de unidad: metodos individuales.
    \item Testing de integracion: metodos de una o varias clases.
    \item Testing de sistema: el programa completo.
  \end{itemize}

  \item \textbf{Listar todas las entradas.}
  \begin{itemize}
    \item Metodos: parametros y estado (variables no locales).
    \item Modulo: parametros y variables globales involucradas.
    \item Sistema: archivos, bases de datos, variables de entorno, etc.
  \end{itemize}

  \item \textbf{Crear un modelo del dominio de las entradas (MDE).}
  \begin{itemize}
    \item Identificar \textbf{caracteristicas} relevantes.
    \item Particionar cada caracteristica en \textbf{bloques}.
    \item Verificar completitud y disjointness.
  \end{itemize}

  \item \textbf{Aplicar un criterio} para generar combinaciones de bloques (ECC, PWC, BCC, ACoC).

  \item \textbf{Obtener entradas concretas} a partir de las combinaciones (valores reales para
  ejecutar los tests).
\end{enumerate}

Los pasos 1 y 2 estan en el nivel de abstraccion \emph{concreto} (codigo real, parametros). Los
pasos 3 y 4 se realizan en el nivel \emph{de diseno} (modelo abstracto). El paso 5 vuelve al nivel
concreto.

\begin{ideabox}
Un MDE \textbf{representa el espacio de entradas de un SUT (System Under Test) en forma
abstracta}. Permite razonar sobre que se prueba sin atarse a valores concretos.
\end{ideabox}

\subsection{Dos enfoques para construir el MDE}

\subsubsection{Enfoque basado en la interfaz}

\begin{itemize}
  \item Se generan caracteristicas a partir de cada parametro \emph{aisladamente}.
  \item Es basado en sintaxis (tipos, presencia de \texttt{null}, etc.).
  \item Suele ignorar informacion semantica y las relaciones entre parametros.
  \item Es facil y mecanico, pero limitado.
\end{itemize}

\subsubsection{Enfoque basado en la funcionalidad}

\begin{itemize}
  \item Las caracteristicas reflejan el \emph{comportamiento esperado} del programa.
  \item Permite incorporar informacion semantica y del dominio.
  \item Captura relaciones entre parametros (por ejemplo, ``\(a+b > c\)'' para un triangulo
  valido).
  \item Requiere mas esfuerzo de diseno.
\end{itemize}

\begin{ideabox}
\textbf{Heuristica de Ammann \& Offutt}: las \emph{precondiciones} y \emph{poscondiciones} son
fuentes excelentes de caracteristicas. Separan explicitamente el comportamiento normal del
excepcional.
\end{ideabox}

\subsubsection{Ejemplo: \texttt{kindOfTriangle(int side1, int side2, int side3)}}

\textbf{Enfoque por interfaz (refinado)}: una caracteristica por lado con bloques
\(\{>1,\ =1,\ =0,\ <0\}\). Hay \(4 \times 4 \times 4 = 64\) combinaciones posibles pero solo
algunas corresponden a triangulos validos.

\textbf{Enfoque por funcionalidad (semantico)}: una caracteristica unica ``clasificacion
geometrica'' con bloques \(\{\text{escaleno}, \text{isosceles no equilatero}, \text{equilatero},
\text{invalido}\}\). Notar que se ajusta el bloque ``isosceles'' a ``isosceles no equilatero''
para que sea \emph{disjunto} de ``equilatero''.

Una alternativa intermedia es usar varias caracteristicas booleanas con \emph{restricciones}:
\begin{itemize}
  \item \(q_1 = \text{escaleno},\ q_2 = \text{isosceles},\ q_3 = \text{equilatero},\ q_4 = \text{valido}\).
  \item Restricciones: \(q_3 = \text{true} \Rightarrow q_2 = \text{true}\);
  \(q_4 = \text{false} \Rightarrow q_1 = q_2 = q_3 = \text{false}\).
\end{itemize}

\subsubsection{Ejemplo: \texttt{numberOfOcurrences(List, element)}}

\begin{itemize}
  \item \textbf{Por interfaz}: dos caracteristicas booleanas, una por parametro (``\texttt{list}
  es null'', ``\texttt{list} es vacia''). No capta cuantas veces aparece el elemento.
  \item \textbf{Por funcionalidad}:
  \begin{itemize}
    \item ``Numero de ocurrencias'' \(\in \{0, 1, >1\}\),
    \item ``\texttt{element} esta primero en la lista'' \(\in \{\text{true},\text{false}\}\),
    \item ``\texttt{element} esta ultimo en la lista'' \(\in \{\text{true},\text{false}\}\).
  \end{itemize}
\end{itemize}

\subsection{Valores de test y estrategias}

Una vez fijadas las particiones, hay que elegir valores reales para cada bloque. Recomendaciones:

\begin{itemize}
  \item Incluir valores que representen \textbf{uso normal}.
  \item Incluir \textbf{valores invalidos} y \textbf{valores especiales} (\texttt{null}, 0, cadenas
  vacias, true/false, espacios en blanco).
  \item Considerar \textbf{valores limite} (borders).
  \item Re-particionar bloques que sean ``muy gruesos''.
  \item Balancear el numero de bloques entre caracteristicas.
  \item Verificar disjointness y completitud antes de generar tests.
\end{itemize}

\begin{ejemplobox}
Si la particion para ``relacion con cero'' es \(\{>1, =1, =0, <0\}\), un buen conjunto de valores
es \(\{2, 1, 0, -1\}\) en lugar de \(\{5, 1, 0, -5\}\), porque incluye los \textbf{bordes}
\(1\) y \(-1\), que son los que disparan defectos tipicos (errores ``off-by-one'').
\end{ejemplobox}

\section{Criterios de cobertura sobre el espacio de entradas}

Una vez definidas las caracteristicas y los bloques, hay que decidir como \emph{combinarlos} para
producir requisitos de test. Los criterios principales son:

\subsection{All Combinations Coverage (ACoC)}

\begin{definicionbox}
Se generan todos los productos cartesianos posibles de bloques entre caracteristicas.

Si hay \(k\) caracteristicas con \(n_1, \ldots, n_k\) bloques cada una, el numero de tests es
\(n_1 \times n_2 \times \cdots \times n_k\).
\end{definicionbox}

Es el criterio mas exhaustivo. En la practica casi nunca se usa porque explota combinatoricamente:
para \texttt{kindOfTriangle} con tres caracteristicas de 4 bloques son \(4 \times 4 \times 4 = 64\)
tests, de los cuales solo unos pocos corresponden a triangulos validos.

\subsection{Each Choice Coverage (ECC)}

\begin{definicionbox}
Cada bloque de cada caracteristica debe aparecer en al menos un test.

El numero de tests es igual al numero de bloques de la \emph{caracteristica con mas bloques}.
\end{definicionbox}

Es el criterio mas barato pero el menos exigente: no garantiza interacciones entre caracteristicas.
Suele detectar pocos defectos.

\begin{ejemplobox}
Para \texttt{kindOfTriangle} con tres caracteristicas de 4 bloques cada una, alcanza con 4 tests:
por ejemplo \((2,2,2),(1,1,1),(0,0,0),(-1,-1,-1)\). Cada bloque de cada caracteristica aparece al
menos una vez.
\end{ejemplobox}

\subsection{Pair-Wise Coverage (PWC) y t-Wise Coverage (TWC)}

\begin{definicionbox}
\textbf{PWC.} Para cada par de bloques de caracteristicas distintas, debe existir al menos un
test que los cubra.

\textbf{TWC.} Generalizacion: cada tupla de \(t\) bloques de \(t\) caracteristicas distintas debe
estar cubierta. PWC es \(t = 2\); ACoC es \(t = k\) (numero total de caracteristicas).
\end{definicionbox}

El numero de tests para PWC es \emph{al menos} el producto de los bloques de las dos
caracteristicas con mas bloques.

\textbf{Por que PWC es popular}: la evidencia empirica sugiere que la mayoria de los defectos en
software se disparan por la interaccion de a lo sumo dos parametros. PWC ofrece una excelente
relacion costo/beneficio respecto a ACoC.

\begin{ejemplobox}
Tres particiones con bloques \([A,B]\), \([1,2,3]\) y \([x,y]\) generan
\(2 \times 3 + 2 \times 2 + 3 \times 2 = 16\) pares. Como cada test cubre 3 pares (uno por cada
combinacion de caracteristicas), alcanzan 8 tests para cubrirlos todos:
\((A,1,x),(A,2,x),(A,3,x),(B,1,y),(B,2,y),(B,3,y),(A,1,y),(B,1,x)\).
\end{ejemplobox}

\subsection{Base Choice Coverage (BCC)}

\begin{definicionbox}
Para cada caracteristica se designa un bloque como \textbf{eleccion base}. El test base se forma
con la eleccion base de todas las caracteristicas. Los demas tests se forman manteniendo todas las
bases constantes excepto una, que toma cada uno de los bloques no base.

Numero de tests: \(1 + \sum_i (n_i - 1)\), donde \(n_i\) es el numero de bloques de la
caracteristica \(i\).
\end{definicionbox}

La intuicion es: \emph{fijar un escenario tipico (camino feliz) y variar una caracteristica a la
vez}. Es muy util cuando hay un caso ``normal'' claramente identificable y se quiere medir el
impacto de cada variacion.

\begin{ideabox}
\textbf{El test base debe ser viable.} Tipicamente conviene elegir el ``camino feliz'' o el caso
mas frecuente como base. La eleccion de las bases es una decision de diseno; conviene
documentarla.
\end{ideabox}

\begin{ejemplobox}
Para tres caracteristicas A, B, C con 4 bloques cada una, eligiendo \((A_1, B_1, C_1)\) como
base, BCC produce \(1 + 3 + 3 + 3 = 10\) tests:
\((A_1,B_1,C_1),(A_2,B_1,C_1),(A_3,B_1,C_1),(A_4,B_1,C_1),(A_1,B_2,C_1),(A_1,B_3,C_1),
(A_1,B_4,C_1),(A_1,B_1,C_2),(A_1,B_1,C_3),(A_1,B_1,C_4)\).
\end{ejemplobox}

Existe una variante llamada \textbf{Multiple Base Choice (MBCC)} en la que se eligen varias bases
distintas y se repite el procedimiento para cada una.

\subsection{Restricciones entre caracteristicas}

No todas las combinaciones de bloques son factibles. Las restricciones son de dos tipos:

\begin{itemize}
  \item Un bloque de una caracteristica \emph{no puede} combinarse con un bloque especifico de
  otra (combinacion prohibida).
  \item Un bloque \emph{solo puede} combinarse con un bloque especifico de otra (combinacion
  forzada).
\end{itemize}

\textbf{Como se manejan segun el criterio}:
\begin{itemize}
  \item \textbf{ACoC y PWC}: se eliminan las combinaciones no viables del conjunto de TRs.
  \item \textbf{BCC}: si la variacion produce una combinacion no viable, se cambia uno de los
  valores base por otro no base hasta lograr una combinacion factible.
\end{itemize}

\begin{ejemplobox}
Para \texttt{numberOfOcurrences} con caracteristicas:
\begin{itemize}
  \item A: ocurrencia del elemento \(\in \{\text{no aparece}, \text{una vez}, \text{mas de una}\}\),
  \item B: primero en la lista \(\in \{\text{true}, \text{false}\}\),
  \item C: ultimo en la lista \(\in \{\text{true}, \text{false}\}\),
\end{itemize}
hay dos restricciones obvias:
\begin{itemize}
  \item \(\neg(\text{no aparece} \wedge \text{primero} = \text{true})\),
  \item \(\neg(\text{no aparece} \wedge \text{ultimo} = \text{true})\).
\end{itemize}
Combinaciones invalidas: \((A.B_1, B.B_1)\) y \((A.B_1, C.B_1)\).
\end{ejemplobox}

\subsection{Comparacion rapida de criterios}

\begin{longtable}{@{}p{2.5cm}p{4cm}p{4cm}p{4cm}@{}}
\toprule
\textbf{Criterio} & \textbf{Idea} & \textbf{\# Tests} & \textbf{Comentario} \\
\midrule
\endhead
ECC & Cada bloque aparece al menos una vez. & \(\max_i n_i\) & Mas barato, menos efectivo. \\
PWC & Cada par de bloques de caracteristicas distintas aparece. & Al menos el producto de las dos caracteristicas con mas bloques. & Buen balance costo/beneficio. \\
BCC & Test base + una variacion por bloque no base. & \(1 + \sum (n_i - 1)\). & Util cuando hay un ``camino feliz'' claro. \\
ACoC & Producto cartesiano completo. & \(\prod n_i\). & Maximo nivel de combinacion, suele ser inviable. \\
\bottomrule
\end{longtable}

\textbf{Relacion de subsumption}: ACoC subsume a PWC, que subsume a ECC. BCC \emph{no} se ordena
linealmente con los anteriores: en general no subsume ni es subsumido por PWC, porque cubre
escenarios distintos (variaciones controladas en lugar de pares forzados).

\section{Conexion con los ejercicios del TP3}

El TP3 instancia estos conceptos sobre cinco unidades concretas:

\begin{itemize}
  \item \textbf{Ejercicio 2 -- PWC sobre \texttt{numberOfOcurrences}.} MDE basado en
  funcionalidad con cuatro caracteristicas (referencias y forma de la lista). Las restricciones
  son centrales para que los TRs sean factibles.

  \item \textbf{Ejercicio 3 -- analisis de un MDE + BCC sobre \texttt{intersection}.} El MDE
  original viola \emph{completitud} y \emph{disjointness} en la caracteristica que describe la
  relacion entre los sets; el ejercicio pide encontrar contraejemplos y proponer un MDE corregido
  antes de aplicar BCC.

  \item \textbf{Ejercicio 4 -- PWC sobre \texttt{PatternIndex}.} MDE combinado (interfaz +
  funcionalidad) con caracteristicas para el resultado esperado. La suite revela dos defectos
  clasicos: precondiciones no validadas y manejo incorrecto de cadenas vacias.

  \item \textbf{Ejercicio 5 -- PWC sobre \texttt{Iterator}.} Caracteristicas booleanas que
  describen el estado del iterador. Como \texttt{Iterator} es una interfaz, hay que elegir
  implementaciones concretas para cada combinacion (por ejemplo, usar
  \texttt{Collections.unmodifiableList} para forzar el caso ``\texttt{remove()} no soportado'').
\end{itemize}

\section{Glosario rapido}

\begin{itemize}[itemsep=0.15em]
  \item \textbf{SUT}: \textit{System Under Test}. La unidad/modulo/sistema que se prueba.
  \item \textbf{MDE}: Modelo del Dominio de Entradas (en ingles, IDM, \textit{Input Domain
  Model}).
  \item \textbf{TR}: Test Requirement, requisito de test.
  \item \textbf{ECC}: \textit{Each Choice Coverage}.
  \item \textbf{PWC}: \textit{Pair-Wise Coverage}; caso particular de TWC con \(t=2\).
  \item \textbf{TWC}: \textit{t-Wise Coverage}; generaliza PWC a tuplas de \(t\) caracteristicas.
  \item \textbf{BCC}: \textit{Base Choice Coverage}.
  \item \textbf{ACoC}: \textit{All Combinations Coverage}.
  \item \textbf{MBCC}: \textit{Multiple Base Choice Coverage}; variante con varias bases.
  \item \textbf{Subsumption}: relacion entre criterios, \(C_1\) subsume a \(C_2\) si toda suite
  que cumple \(C_1\) cumple \(C_2\).
  \item \textbf{Infeasible TR}: requisito de test que ningun input puede satisfacer.
\end{itemize}

\section*{Cierre}
\addcontentsline{toc}{section}{Cierre}

El particionado del espacio de entradas es una de las herramientas mas versatiles de testing
moderno: \emph{no requiere conocer la implementacion}, solo la especificacion del espacio de
entradas; se adapta a unit, integration y system testing; y permite ajustar el esfuerzo cambiando
de criterio. La diferencia entre una suite ``casual'' y una basada en criterios es exactamente la
trazabilidad: cada test responde a un TR, y la suite completa se puede auditar y mantener con
claridad.

\end{document}
